% Options for packages loaded elsewhere
\PassOptionsToPackage{unicode}{hyperref}
\PassOptionsToPackage{hyphens}{url}
\documentclass[
]{article}
\usepackage{xcolor}
\usepackage[margin=1in]{geometry}
\usepackage{amsmath,amssymb}
\setcounter{secnumdepth}{-\maxdimen} % remove section numbering
\usepackage{iftex}
\ifPDFTeX
  \usepackage[T1]{fontenc}
  \usepackage[utf8]{inputenc}
  \usepackage{textcomp} % provide euro and other symbols
\else % if luatex or xetex
  \usepackage{unicode-math} % this also loads fontspec
  \defaultfontfeatures{Scale=MatchLowercase}
  \defaultfontfeatures[\rmfamily]{Ligatures=TeX,Scale=1}
\fi
\usepackage{lmodern}
\ifPDFTeX\else
  % xetex/luatex font selection
\fi
% Use upquote if available, for straight quotes in verbatim environments
\IfFileExists{upquote.sty}{\usepackage{upquote}}{}
\IfFileExists{microtype.sty}{% use microtype if available
  \usepackage[]{microtype}
  \UseMicrotypeSet[protrusion]{basicmath} % disable protrusion for tt fonts
}{}
\makeatletter
\@ifundefined{KOMAClassName}{% if non-KOMA class
  \IfFileExists{parskip.sty}{%
    \usepackage{parskip}
  }{% else
    \setlength{\parindent}{0pt}
    \setlength{\parskip}{6pt plus 2pt minus 1pt}}
}{% if KOMA class
  \KOMAoptions{parskip=half}}
\makeatother
\usepackage{graphicx}
\makeatletter
\newsavebox\pandoc@box
\newcommand*\pandocbounded[1]{% scales image to fit in text height/width
  \sbox\pandoc@box{#1}%
  \Gscale@div\@tempa{\textheight}{\dimexpr\ht\pandoc@box+\dp\pandoc@box\relax}%
  \Gscale@div\@tempb{\linewidth}{\wd\pandoc@box}%
  \ifdim\@tempb\p@<\@tempa\p@\let\@tempa\@tempb\fi% select the smaller of both
  \ifdim\@tempa\p@<\p@\scalebox{\@tempa}{\usebox\pandoc@box}%
  \else\usebox{\pandoc@box}%
  \fi%
}
% Set default figure placement to htbp
\def\fps@figure{htbp}
\makeatother
\setlength{\emergencystretch}{3em} % prevent overfull lines
\providecommand{\tightlist}{%
  \setlength{\itemsep}{0pt}\setlength{\parskip}{0pt}}
\usepackage{bookmark}
\IfFileExists{xurl.sty}{\usepackage{xurl}}{} % add URL line breaks if available
\urlstyle{same}
\hypersetup{
  pdftitle={Cognitive Models Research Report},
  pdfauthor={{[}YOUR NAME},
  hidelinks,
  pdfcreator={LaTeX via pandoc}}

\title{Cognitive Models Research Report}
\usepackage{etoolbox}
\makeatletter
\providecommand{\subtitle}[1]{% add subtitle to \maketitle
  \apptocmd{\@title}{\par {\large #1 \par}}{}{}
}
\makeatother
\subtitle{An Investigation into the StarWars Data}
\author{{[}YOUR NAME}
\date{2025-11-17}

\begin{document}
\maketitle

{
\setcounter{tocdepth}{2}
\tableofcontents
}
The Model: Hyperbolic Temporal DiscountingPsychological Relevance: This
is the quintessential model of impulsivity, self-control, and addiction.
It's used to explain why people (or animals) often choose a small,
immediate reward (e.g., drinking, procrastination) over a large, delayed
reward (e.g., health, a finished project). It bridges the gap between
learning (R-W) and decision-making.The Assignment Idea:Implement the
Model: Have students write a function for the hyperbolic discounting
equation: \(V = A / (1 + kD)\), where \(V\) is subjective value, \(A\)
is the reward amount, \(D\) is the delay, and \(k\) is the individual's
discount rate (impulsivity).Simulate a ``Preference Reversal'': Ask them
to model a choice between a Small-Immediate Reward (e.g., \(A=10, D=0\))
and a Large-Delayed Reward (e.g., \(A=20, D=10\)).Run the Simulation:
Have them plot the subjective value (\(V\)) of both options as the delay
to the Small-Immediate reward counts down (e.g., from \(D=20\) to
\(D=0\)).Key Insight (The Graph): The student's plot should show a
``crossover.'' Far away, the large reward is always more valuable. But
as the small reward becomes immediately available, its value spikes,
causing the ``preference reversal'' that models a lapse in
self-control.Connect to Psychology: Ask them to plot two lines: one for
a ``low \(k\)'' individual (patient) and one for a ``high \(k\)''
individual (impulsive) and explain the difference.

The Model: The ``Multi-Armed Bandit'' (Q-Learning) 🎰 Psychological
Relevance: This is a simple but powerful Reinforcement Learning (RL)
model. It's the most natural next step from Rescorla-Wagner. While R-W
learns to predict an outcome (``the bell means food is coming''), a
Q-learning model learns to make a choice to get a reward (``I should
pull this lever to get food''). It's the classic model for the
exploration vs.~exploitation trade-off, which is fundamental to all
decision-making.

The Assignment Idea:

Set up the ``Casino'': Create a simulation with three ``slot machines''
(three ``armed bandits''). Each machine has a different, hidden
probability of paying out (e.g., Machine A pays 50\% of the time, B pays
30\%, C pays 70\%).

Implement the ``Agent'': The student creates an ``agent'' that starts
with no knowledge. The agent's ``brain'' is a simple vector of Q-values
(Quality values), one for each machine, all starting at 0.

The ``Learning'' Loop (Trial-by-Trial):

The agent has to choose a machine to pull. This is where you introduce
the Softmax Decision Rule. The agent is more likely to pick machines
with high Q-values, but there's a small chance (regulated by a
``temperature'' parameter) it will explore a low-Q-value machine.

After pulling, the ``casino'' (your code) gives a reward (1 or 0).

The agent then updates its Q-value for only the machine it chose, using
an update rule that looks just like Rescorla-Wagner: Q{[}chosen{]} =
Q{[}chosen{]} + alpha * (reward - Q{[}chosen{]})

Key Insight (The Graph): Ask them to plot two things: 1) The Q-values
over time (they should converge on the true probabilities) and 2) The
percent of choices for each machine over time (they should see the agent
``discover'' the 70\% machine and eventually exploit it).

The Model: Schelling's Segregation Model (Agent-Based) 🏘️ Psychological
Relevance: This is a classic agent-based model (ABM) from social
psychology/sociology. It's a powerful demonstration of how complex,
large-scale social patterns (like segregated neighborhoods) can emerge
from very simple, seemingly harmless individual ``rules'' or
preferences.

The Assignment Idea:

Create the ``City'': Have the student create an N x N matrix (e.g.,
20x20) to represent the city.

Populate the City: Randomly sprinkle two ``types'' of agents (e.g., A
and B) into the city, leaving some empty spots.

Define the ``Happiness'' Rule (The Cognitive Part): Each agent has a
simple preference: ``I will be `happy' as long as at least X\% of my 8
neighbors are like me.'' The key is to set X to a low, non-racist value,
like 30\% (i.e., ``I don't want to be the only one'').

The ``Simulation'' Loop:

Write a loop that runs for 50 ``rounds.''

In each round, the code finds all ``unhappy'' agents.

It then moves each unhappy agent to a random empty spot in the city.

Key Insight (The Graph): Ask them to plot the state of the city (the
matrix) at Round 1 (random chaos) and at Round 50. They will be shocked
to see that the city has become highly segregated. The ``A''s will all
be in one part of town, and the ``B''s in another, even though no agent
``wanted'' segregation. This is a profound lesson in emergent social
phenomena.

\section{Decision-making}\label{decision-making}

The Model: A ``Two-Alpha'' Reinforcement Learning AgentInstead of one
``association'' (\(\Delta V\)) like in R-W, this agent tracks an
Expectancy (\(E\)) for each of the four decks (A, B, C, D), all starting
at 0.The key insight is that healthy people and patients learn from
gains and losses differently. We can model this with two separate
learning rates:\(\alpha_{gain}\) (Alpha-Gain): How much the agent learns
from a better-than-expected outcome.\(\alpha_{loss}\) (Alpha-Loss): How
much the agent learns from a worse-than-expected outcome.The Update Rule
(The ``Delta Rule''):Choose a deck (e.g., Deck A).Get the outcome (e.g.,
Win \$100, Lose \$0).Calculate the Prediction Error (PE): PE = (Win -
Loss) - E{[}Deck A{]}If PE \textgreater{} 0 (it's a gain), update using
\(\alpha_{gain}\):E{[}Deck A{]} = E{[}Deck A{]} + alpha\_gain * PEIf PE
\textless{} 0 (it's a loss), update using \(\alpha_{loss}\):E{[}Deck
A{]} = E{[}Deck A{]} + alpha\_loss * PEThe ``Cognitive Knob'' (The
Psychology):Healthy Agent: Symmetrical learning. alpha\_gain = 0.1 and
alpha\_loss = 0.1. They learn from pain just as much as they learn from
pleasure.''Patient'' Agent (Myopic): Asymmetrical learning. alpha\_gain
= 0.1 but alpha\_loss = 0.01. They are excited by gains but functionally
ignore or underweight the long-term punishments.The Assignment: Modeling
the Iowa Gambling TaskObjective: Simulate a ``Healthy'' and a
``Patient'' agent to understand how a simple cognitive bias (ignoring
losses) leads to self-defeating choices.Part 1: Build the ``Casino''
(The IGT Environment)First, students must write a function that
simulates the IGT. Give them this payoff schedule:DeckRewardLoss
FrequencyLoss
AmountA+\$10050\%-\$250B+\$10010\%-\$1250C+\$5050\%-\$50D+\$5010\%-\$250Net
(A/B): -\$25 per 10 cards. Net (C/D): +\$25 per 10 cards.Task: Write a
function pull\_card(deck) that, when called, returns a list(win, loss)
for one of the four decks.Part 2: Build the ``Agent'' (The Cognitive
Model)Task: Write a function update\_expectancy(E\_old, win, loss,
alpha\_gain, alpha\_loss) that takes the old expectancy and the outcome,
and returns the new expectancy using the two-alpha delta rule.Part 3:
The Choice Rule (Softmax)How does the agent choose? It uses the Softmax
(or Boltzmann) decision rule. This rule translates the expectancy values
into choice probabilities.Task: Give them the Softmax equation and have
them write a function choose\_deck(E\_vector, tau).E\_vector is the list
of expectancies (e.g., c(E\_A, E\_B, E\_C, E\_D)).tau (Temperature) is a
``consistency'' parameter.1 A low tau means the agent always picks the
highest E (exploits). A high tau means it chooses randomly
(explores).Part 4: The Simulation (Putting it all together)Task: Run a
simulation for two agents (Healthy vs.~Patient) over 100 trials.Healthy
Agent Simulation:Set alpha\_gain = 0.1, alpha\_loss = 0.1.Loop 100
times. In each loop:Choose a deck using choose\_deck().Get the (win,
loss) from your pull\_card() function.Update the expectancy for that
deck only using update\_expectancy().Record which deck was
chosen.Patient Agent Simulation:Do the exact same thing, but with
alpha\_gain = 0.1 and alpha\_loss = 0.01.Part 5: The ``Aha!'' Moment
(The Graph)Task: Plot the results. The best way is to plot the
percentage of choices from ``Bad Decks'' (A \& B) over time.Create a
plot showing ``Trial Number'' (1-100) on the x-axis.On the y-axis, show
the cumulative percentage of choices made from Decks A \& B.Plot two
lines: one for the Healthy Agent, one for the Patient Agent.Expected
Insight:The Healthy Agent's line will decrease over time. They will
start by sampling all decks, but as they learn from the large
punishments, they will ``learn'' to prefer the ``Good'' decks C \& D.The
Patient Agent's line will stay high or even increase. Because they
ignore the losses (alpha\_loss = 0.01), the big +\$100 rewards from
Decks A \& B look amazing to them. They never learn that these decks are
bankrupting them.This directly models the myopic, self-defeating
behavior seen in the task and provides a clear cognitive mechanism
(asymmetrical learning) for why it happens.

\end{document}
